%=====================================================================
\clearpage
\hdr{A}{附录 A：想深一层（不读也能做完实验）}

这三节回答「为什么这个仪器这样做」，与操作无关。现场时间紧就跳过，报告写完再看。

\subhdr{A.1 图 12-1 的 28 个编号：你真正会动的只有 12 个}

不必全记。按「什么时候动它」分三组，编号与原名不变：

\begin{center}
\begin{minipage}{\linewidth}
\small
\begin{tabular}{@{}p{0.6cm}p{6.0cm}p{7.4cm}@{}}
\toprule
\thead{编号} & \thead{名称（英文原名）} & \thead{什么时候动它} \\
\midrule
\multicolumn{3}{@{}l}{\textbf{A. 望远镜这一路}}\\
8 & 望远镜 \emph{telescope} & 转它去找像、找谱线 \\
9 & 目镜筒锁紧 \emph{locking screw of eyepiece} & 松开后目镜筒才能前后移 \\
11 & 目镜视度调节手轮 \emph{adjustment hand wheel} & \textbf{最先动它}：把叉丝本身看清 \\
12 & 望远镜光轴竖直调节 \emph{vertical adjustment screw} & \textbf{各半法调它} \\
13 & 望远镜光轴水平调节 \emph{horizontal adjustment screw} & 目测粗调时对准准直管 \\
15 & 望远镜微调 \emph{fine tuning screw} & 把像对到叉丝中心 \\
17 & 望远镜锁紧 \emph{locking screw of telescope} & 对准后锁死 \\
\midrule
\multicolumn{3}{@{}l}{\textbf{B. 准直管这一路}}\\
1 & 狭缝装置 \emph{slit device} & 转它让狭缝像竖直 \\
2 & 狭缝锁紧 \emph{slit device locking screw} & 松—移—紧一次完成，之后别碰 \\
26 & 准直管光轴水平调节 \emph{horizontal screw} & 目测粗调时对准望远镜 \\
27 & 准直管光轴竖直调节 \emph{vertical screw} & 步骤 5 调它：像落在视场中央 \\
28 & 缝宽调节手轮 \emph{hand wheel of slit width} & 调到约 $0.5$ mm（范围 $0.2$--$2$ mm） \\
\midrule
\multicolumn{3}{@{}l}{\textbf{C. 载物台、度盘与游标}}\\
5 & 载物台 \emph{specimen stage} & 放平面镜、三棱镜、光栅 \\
6 & 载物台调平螺钉 \emph{leveling bolt}（三个 $a_1a_2a_3$） & \textbf{各半法调它}；三螺钉成正三角形 \\
16 & 游标盘与度盘的锁紧 \emph{locking screw}（背面） & \textbf{决定「谁跟谁一起转」}；一组测量开始后不许改 \\
21 & 游标 \emph{vernier}（两个，相差 $180^{\circ}$） & \textbf{读数看它} \\
22 & 度盘 \emph{dial plate}（$0$--$360^{\circ}$） & 读数看它 \\
24 & 游标盘微调 \emph{fine tuning screw of alidade} & 把像对到叉丝中心 \\
25 & 游标盘锁紧 \emph{locking screw of alidade} & 对准后锁死 \\
7 & 载物台锁紧 & 找到最小偏向角后锁台 \\
\bottomrule
\end{tabular}
\end{minipage}
\end{center}

%=====================================================================
\needspace{16\baselineskip}\hdr{B}{附录 B：两份材料说法不一致时，以这条为准}

讲义（15 页）与预习（3 页）有八处彼此不一致或没写清楚。现场遇到就照下表，不要在仪器前纠结。

\begin{center}
\begin{minipage}{\linewidth}
\small
\begin{tabular}{@{}p{4.0cm}p{5.2cm}p{5.0cm}@{}}
\toprule
\thead{问题} & \thead{材料怎么说} & \thead{本讲口径} \\
\midrule
重复几次 & 正文说 five times，预习表只留 3 行 & 按 3 组做；老师要 5 次就加两行 \\
光栅线数 & 只给 $N=1/(1000d)$，没给具体数 & 先看光栅框上的标注；未标注按 $300$ 线/mm 预估角度 \\
棱镜是否等边 & 未写 & 按等边（$\alpha=60^{\circ}$）做例算与判据 \\
游标「取零线右边」 & 与它自己的算例矛盾 & 以算例为准：取零线\textbf{已越过}的那条 \\
$(12\text{-}1)(12\text{-}2)$ & 按 $\angle1=\angle2$ 代入会得 $\Phi=\theta_1$，与结论矛盾 & 属笔误；结论 $\phi=(\theta_1+\theta_2)/2$ 不变（附录 A.2 重推） \\
\emph{dispersive power} & 引言用这个词，正文定义的却是 $D=\Delta\varphi/\Delta\lambda$ & 一律按\textbf{角色散}处理 \\
图 12-12 角标 & 第二个游标印成 $\phi_2''$ & 统一用 $\phi_2'$ \\
两套螺钉编号 & 图 12-10 用 $a,b,c$，图 12-11 用 $a_1,a_2,a_3$ & 不是同一套；现场先看载物台刻字，按「只改一个面」的原则挑 \\
钠灯 & 器材清单列了，流程全用汞灯 & 本讲不给钠灯操作，不编造步骤 \\
\bottomrule
\end{tabular}
\end{minipage}
\end{center}

%=====================================================================
\needspace{16\baselineskip}\hdr{C}{附录 C：最贵的十个错}

\begin{center}
\begin{minipage}{\linewidth}
\small
\begin{tabular}{@{}p{0.6cm}p{5.6cm}p{7.6cm}@{}}
\toprule
\thead{\#} & \thead{错法} & \thead{后果与正确做法} \\
\midrule
1 & 只读一个游标 & 引入最多 $1'$ 的偏心误差。\textbf{每读一次方位必须两个游标都读} \\
2 & 「像找到了」当成「调好了」 & 绿十字模糊或带视差就开始测，所有角量带固定偏差。判据是\textbf{清晰 + 无视差 + 与上叉丝重合}三件齐 \\
3 & 各半法一次调到位 & 一面重合、另一面偏差更大且不收敛。正确是\textbf{先减半、再归零} \\
4 & 折回 $360^{\circ}$ 只在平均后做一次 & 某行原始差 $240^{\circ}$ 当 $240^{\circ}$ 参与平均，$\alpha$ 错几十度。\textbf{每行都判断} \\
5 & $d$ 与 $N$ 混用 / 单位没换 & 结果差 $10^3$ 量级。\textbf{先把 $d$ 换到 nm} \\
6 & 调 $a_2$ 想「一起摆平」 & 两个面被同时扰动。\textbf{只调 $a_1,a_3$} \\
7 & 测最小偏向角时转太快 & 跳过折返点，$\delta$ 偏大。\textbf{慢转、盯住停顿} \\
8 & 光栅只调 $b_1,b_2$ 没调 $b_3$ & 谱线不等高，叉丝与斜谱线交点含糊 \\
9 & 中途动了 16 号螺钉 & 前后读数不属同一度盘，整组作废。\textbf{一组开始后别碰} \\
10 & 读数写成 $335^{\circ}47'$ 而实际 $334^{\circ}47'$ & 没做「结果落在两刻线之间」的自检 \\
\bottomrule
\end{tabular}
\end{minipage}
\end{center}

\warnbox{\textbf{时间不够时只保三件事}：① 望远镜聚焦到无视差；② 望远镜轴垂直主轴（镜面两面都重合）；③ 每次读数都读两个游标。这三件保证数据的系统性；其余的（少测一组、二级没测）只是精度损失，可以事后说明。反过来，这三件缺任何一件，数据不可救。}

\subhdr{A.2 自准直为什么比肉眼准：灵敏度放大一倍}

设平面镜相对光轴偏了 $\gamma$，由反射定律，返回光的方向偏了 $2\gamma$，像点在上叉丝旁移开

\[
s=2f\gamma .
\]

取 $f=200$ mm、$\gamma=1'=2.91\times10^{-4}$ rad，则 $s=2\times200\times2.91\times10^{-4}=0.116$ mm——肉眼清楚可见。

这就是自准直的价值：\textbf{角度的判据被放大了一倍，再被焦距换成一段可以直接看的位移}。$1'$ 的偏角在视场里是一段看得见的距离，所以「像与上叉丝重合」比「看上去大致对准」精确得多。反过来说，$f$ 越大越灵敏，但视场越小、越难找像——这就是为什么粗调必须先把像找进视场。

\subhdr{A.3 偏心差的一阶推导（为什么两个游标取平均就消掉）}

设主轴中心 $O$ 在原点，度盘中心 $O'$ 在 $(e\cos\psi,\ e\sin\psi)$，度盘半径 $R$。游标 $V$ 随望远镜绕 $O$ 转到真实方位角 $\phi$，而度盘刻线从 $O'$ 读出角 $\theta$。于是

\[
V=\bigl(e\cos\psi+R\cos\theta,\ e\sin\psi+R\sin\theta\bigr).
\]

从 $O$ 看 $V$ 的方向满足

\[
\tan\phi=\frac{R\sin\theta+e\sin\psi}{R\cos\theta+e\cos\psi}
=\tan\theta+\frac{(e/R)\sin(\psi-\theta)}{\cos^2\theta}+\cdots
\]

一阶近似给出 $\phi-\theta\approx(e/R)\sin(\psi-\theta)$，即读数误差

\[
\varepsilon(\theta)=\theta-\phi\approx\frac{e}{R}\sin(\theta-\psi).
\]

两个游标装在一条直径的两端，方位差 $180^{\circ}$：

\[
\varepsilon(\theta+180^{\circ})=\frac{e}{R}\sin(\theta+180^{\circ}-\psi)=-\frac{e}{R}\sin(\theta-\psi)=-\varepsilon(\theta).
\]

\textbf{等大反号}，取平均即消：

\[
\theta_1=\phi+\varepsilon,\quad\theta_2=\phi-\varepsilon
\quad\Longrightarrow\quad
\phi=\frac{\theta_1+\theta_2}{2}.
\]

副产品：$\varepsilon=(\theta_1-\theta_2)/2$，也就是\textbf{两个游标的读数差本身量出了偏心误差}——现场可用它做自检（差应在几个角分内）。

\textbf{量级}：$e=0.02$ mm、$R=60$ mm 时，$\varepsilon_{\max}=\arctan(0.02/60)=1.15'\approx1'$，比读数精度（$0.5'$）还大；平均后只剩 $(e/R)^2$ 量级，约 $0.01''$，完全可忽略。

